\section{Learning a physical correction along the flow}\label{app:force-correction}
The earlier FlowMol study trains a physical head on predicted endpoints with the backbone fixed. The final EGNN uses the same pair construction with a smaller element vocabulary (Appendix~\ref{app:matched-connection}).

\paragraph{States and targets.}
The parent generates two trajectories for each of 128 training compositions, with 32 in each of the 8--16, 17--24, 25--32, and 33--40 atom ranges. We record the second field evaluation at midpoint steps 23 and 29, or times 0.734375 and 0.921875. The resulting 512 states define endpoints $H_t=X_t+(1-t)v_{\rm base}(X_t,t;c)$ after both self-conditioning passes. These are instantaneous endpoint estimates rather than completed trajectories.

Mirror-averaged, centered eSEN forces at $H_t$ give the displacement in Equation~\ref{eq:local-force-shift}, with width at most 0.03\,\AA\ and molecular displacement norm at most 0.1\,\AA. The velocity target is $\delta(H_t)/(1-t)$. We retain all states, including graph-invalid endpoints. Metadata hashes split the parent-training compositions into 96 for fitting and 32 for internal validation, with four states per composition; generator-confirmation compositions are separate.

\paragraph{Equivariant pair network.}
The small head has 8,178 parameters. Symmetric pair features combine atom embeddings, distances in $X_t$ and $H_t$, directional alignment, time, atom count, and covalent radii. Two width-64 SiLU layers predict $\alpha_{ij}$ and $\beta_{ij}$, each bounded to $[-1,1]$ by a hyperbolic tangent.

Let $r_{ij}=X_{t,i}-X_{t,j}$ and $\widehat r_{ij}=H_{t,i}-H_{t,j}$. Define $\bar r_{ij}=r_{ij}/\sqrt{\|r_{ij}\|^2+\ell_{ij}^2}$ and similarly $\overline{\widehat r}_{ij}$, with covalent-radius sum $\ell_{ij}$. Contact weights are $w_{ij}=\operatorname{sigmoid}((1.25-\|\widehat r_{ij}\|/\ell_{ij})/0.2)$. With $d_i=\sum_{j\ne i}w_{ij}$ and $d_{\max}=\max(1,\max_i d_i)$, the correction is
\begin{equation}
 A_{\phi,i}=v_{\rm cap}\,t^2\sum_{j\ne i}
 \frac{w_{ij}}{2d_{\max}}
 \left(\alpha_{ij}\bar r_{ij}+\beta_{ij}\overline{\widehat r}_{ij}\right),
 \qquad v_{\rm cap}=2\text{ \AA\ per unit flow time}.
\end{equation}
Symmetric coefficients and opposite pair vectors give zero total correction and rotation/permutation equivariance \citep{satorras2021egnn}. The triangle inequality gives $\|A_{\phi,i}\|\leq v_{\rm cap}t^2$, a bound on instantaneous velocity addition. Its cumulative effect also depends on the parent flow. Zero initialization of the final layer reproduces the parent before fitting.

Two seeds train for 20,000 AdamW updates using mean squared velocity error, learning rate $3\times10^{-4}$, batch size one, and gradient clipping at one. We also test a 37,540-parameter wider pair network and a 37,586-parameter contextual network with two message-passing rounds and invariant angular summaries. All share the state sequence and targets within each seed.

At 20,000 updates, the contextual head lowers internal-validation target error by approximately 15\% relative to the small head and 11\% relative to the wide head. The common-panel study tests whether this translates to generation quality. Physical residual correction has also been studied for video flows \citep{wang2026pila}; our supervision uses molecular forces at predicted coordinates.

\paragraph{Sampling and preparation.}
We integrate $v_{\rm base}+A_\phi$ for 32 midpoint steps, using 128 backbone and 64 head calls per output without terminal noise or physical queries. A zero-output head matches the parent control's operations. Collecting 512 endpoints uses 256 trajectories and 1,024 mirror-paired eSEN force queries. The shared experiment also evaluates eight translated Gaussian particles per state for a complete-work control, bringing preparation to 9,216 queries; force-target heads do not use those particle energies. Median local-work ESS is 1.37 of eight, and the earlier 2,000-update student comparison found no resolved additional generation gain.
